\section{Encoded cost and the remaining global complexity gap}
\label{sec:complexity}

The size parameters must distinguish four different objects: the knot
diagram, an explicit triangulation at one stage, a binary normal vector,
and a supplied proof.  Let $n$ be the crossing count of a diagram, $t$ the
number of tetrahedra in a particular source triangulation, and
\[
 b=\max\bigl(1,\max_{i,j}\lceil\log_2(1+x_{ij})\rceil\bigr)
\]
the maximum coordinate bit length.  The explicit coordinate input has
$7t$ entries, so reading it already costs $\Omega(tb)$ bits in the worst
case.  Neither $t$ nor $b$ is silently replaced by $n$ during a hierarchy.

For the finite relative compiler, let $s$ count the explicit core handles,
source attaching-cycle occurrences, ports, base-profile records, and
ordered attachment letters.  Let $B\geq\log_2(s+2)$ bound the bit lengths
of all source labels, profile integers, and resulting counters, and let
$M(B)$ denote a multiplication cost for $B$-bit integers.  The number of
private one-handles is a \emph{value} stored in one of these integers; it
is not included in $s$ as a list of handle objects.

\subsection{What the finite compiler actually bounds}

\begin{proposition}[Finite compressed compilation cost]
\label{prop:finite-cost}
Given authenticated finite source geometry, base profiles, and an ordered
attachment skeleton of size $(s,B)$ satisfying the hypotheses of
Theorem~\ref{geom:thm-residual-incidence}, the active incidence matrix has
at most one additional row per attached base and no additional column.
It can be constructed and independently replayed in time polynomial in
$s$ and $B$, independently of the numerical private-handle count and of
the numerical multiplicity of an isolated bundle.

For the delivered straightforward implementation, a conservative bound
is $O(s^2M(B))$ bit work under bounded-cost dictionary indexing of the
explicit labels, with $O(s)$ stored finite records and their $B$-bit
payloads.  Deterministic sorting or balanced dictionaries preserve a
polynomial bound.  The square factor is an implementation bound for
repeated component scans, not an inherent lower bound.
\end{proposition}

\begin{proof}
An original core attaching cycle is scanned once to form its row after
deleting forest columns.  Every source port occurs once in the ordered
attachment components.  A saturated row is the capped sum of the core-one
ports in its own circles.  Since the profiles partition ports, these scans
cannot create a list proportional to the base genus.  The active nonzero
count is bounded by the number of original cycle occurrences plus the
number of core-one ports.

Union--find on the explicit core and attachment records determines which
original cores are joined.  The delivered code then scans profiles,
attachments, handles, and rows for each joined component.  There are at
most $s$ components and $s$ records per scan.  The genus and Euler formulas
use sums and products of encoded integers; charging each such operation
by $M(B)$ gives the stated conservative bound.  The checker reconstructs
the same mathematical quantities with separate code and the same type of
finite scans.  Isolated multiplicity contributes a multiplication to the
histogram rather than repeated graph vertices.  At no point is a private
handle enumerated.
\end{proof}

In the maintained normal-cut source model, core handles and total source
incidences are $O(t)$, and the source port bound is $66t$
\cite{marked-parallelity-85}.  In particular the active matrix has at most
$12t+66t$ rows and $8t$ original columns, using conservative handle and
ported-base bounds.  These are upper bounds, not typical dimensions.
The inherited orbit computations and their encoded output sizes remain
separate costs.  Agol--Hass--Thurston orbit counting is the foundational
compressed connectivity tool\cite{AHT2006}; calling a finite adapter
polynomial does not supply a new proof for every upstream discovery path.

For the matrix layer itself, the stronger bound from
Section~\ref{sec:column-peeling} is available.  With $m$ rows, $k$ active
columns, and $I$ source nonzeros, discovery takes $O(m+k+I)$ expected
dictionary operations, apart from final label sorting.  The counts used
for leaf detection have $O(\log(m+1))$ bits after individual capping.  A
literal count of $2^{16384}$ costs time to read, but never causes that many
updates.

The full census producer has the form
\[
 D_{\rm source}(t,b)+D_{\rm boundary}(t,b)
 +\operatorname{poly}(s,B),
\]
where the first two terms are actual upstream discovery costs.  Its replay
has the analogous source-verification terms plus the cost to read the
supplied orbit traces.  A trace-operation cap bounds the number of trace
records admitted; it is not a hard bound on time, bytes, or malformed
input nesting.  The optional precheck can omit the boundary term on a
certified original-core obstruction.  Successful cases still pay for the
complete source and attachment information.

\subsection{Existence of a large handle system versus outputting it}

Theorem~\ref{geom:thm-census} certifies an invariant of a genuine finite
manifold without printing every disc in a complete disc system.  That
distinction is unavoidable when a profile permits $2^B$ private handles.
A certificate consisting of their count and the residual-incidence
theorem can be short.  A literal output with one named cocore per handle
has an immediate $\Omega(2^B)$ output-size lower bound.

This is why the present compiler cannot simply be declared a general
hierarchy-transition routine.  A later surface might need to cross, or
refer to, individual private cocores.  A short \emph{addressing and
transport scheme} for those cocores is then an additional obligation.
The current proof neither allocates those objects nor gives source
addresses for all of them.  The arbitrary-genus example proves that the
terminal census can remain compressed; it does not bound the cost of
all possible later geometric queries.

\subsection{A conditional composition theorem}

The following elementary theorem makes the global target falsifiable.
Its hypotheses are the missing interface bounds, rather than conclusions
drawn from the local experiment.

\begin{theorem}[Sufficient bounds for a quasi-polynomial composition]
\label{thm:conditional-composition}
Consider a correct recognition algorithm on $n$-crossing diagrams with
at most
\[
 H(n)=L(n)(g(n)+1)^{L(n)}
\]
charged control steps.  Suppose every step tries at most $C(n)$ candidates,
each candidate and its certificate have encoded size at most $S(n)$,
and discovery, transitions, and replay together take at most a fixed
polynomial in $n+C(n)+S(n)$ per charged step.  If
\[
 C(n),S(n),g(n)\leq n^{O(1)},
 \qquad L(n)=O((\log n)^a),
\]
for a fixed $a\geq1$, then its total time is
\[
 2^{O((\log n)^{a+1})}.
\]
In particular $a=1$ suffices for the February slide exponent, and $a=2$
suffices for the March slide exponent.
\end{theorem}

\begin{proof}
Take logarithms of the step bound:
\[
 \log H=\log L+L\log(g+1)
       =O((\log n)^{a+1}).
\]
Multiplying by a fixed polynomial in the stipulated polynomial quantities
adds only $O(\log n)$ to the logarithm of the running time.  A polynomial
preprocessing or final verification phase is absorbed in the same bound.
\end{proof}

The hypothesis includes \emph{finding} the candidate surfaces and
compression discs.  A small verified certificate proves efficient checking
of that certificate; it does not bound a search that might inspect
exponentially many alternatives.  Similarly, a polynomial formula in $t$
and $b$ does not help globally if successive stages make either quantity
exponential in the original $n$.

Even granting $g(n)=n^{O(1)}$, the weaker bound $L(n)=O(n)$ gives
$2^{O(n\log n)}$ by this calculation.  Reducing a local quadratic pass to
a linear pass changes the polynomial factor.  It does not change this
reset exponent.  To improve that exponent, one must either shorten the
hierarchy, reduce the number of resets, or amortize several charged steps
by a stronger theorem.  The ambient-block result later in this article
explores the third possibility under explicit extra geometric hypotheses.

\subsection{Why simple potential shortcuts do not close the gap}

Replacing a lexicographic vector by its sum requires a proof that a repair
does not recreate large later entries.  The sum has no monotonicity merely
because the first altered coordinate decreases.  Likewise, replacing an
integer intersection count $I$ by $\log(I+1)$ does not bound a process
that can delete only one intersection at a time: the logarithmic decrement
can be arbitrarily small.  The potential theorem below retains an honest
numerical intersection term and gives a concrete binary-size obstruction.

These observations guide implementation priorities.  Fast local replay is
valuable and directly measurable.  The global objective additionally needs
a geometrically justified amortization argument, with bounds expressed in
the original input size and maintained throughout every transformation.
